% generated by tools/refs.py from ledger/refs.yaml; do not edit
\begin{thebibliography}{BLGGT14}
\bibitem[AC89]{ArthurClozel1989} \textsc{James Arthur and Laurent Clozel}, \emph{Simple Algebras, Base Change, and the Advanced Theory of the Trace Formula}, Ann. of Math. Stud. 120, Princeton University Press, 1989. \href{https://doi.org/10.1515/9781400882403}{doi:10.1515/9781400882403}
\bibitem[Alp26]{Alpoge2026} \textsc{Levent Alp\"{o}ge}, Remark, note announcing a proof of the Riemann Hypothesis and of GRH, with a Lean archive, 2026. \url{https://alpo.ge/remark.pdf}
\bibitem[Ant26]{GRHLean2026} \textsc{Anthropic, PBC}, grh-lean: the Generalized Riemann Hypothesis for Dirichlet L-functions, in Lean 4, Lean 4 development after the proof of L. Alp\"{o}ge, Apache License 2.0; commit and archive SHA-256 in Appendix E, 2026. \url{https://github.com/lsernau/grh-lean}
\bibitem[Ax70]{Ax1970} \textsc{James Ax}, Zeros of polynomials over local fields---the Galois action, \emph{J. Algebra} \textbf{15} (1970), 417--428. \href{https://doi.org/10.1016/0021-8693(70)90069-4}{doi:10.1016/0021-8693(70)90069-4}
\bibitem[Ber04]{Berger2004} \textsc{Laurent Berger}, Limites de repr\'{e}sentations cristallines, \emph{Compos. Math.} \textbf{140} (2004), 1473--1498. \href{https://doi.org/10.1112/S0010437X04000879}{doi:10.1112/S0010437X04000879}
\bibitem[BG06]{BombieriGubler2006} \textsc{Enrico Bombieri and Walter Gubler}, \emph{Heights in Diophantine Geometry}, New Math. Monogr. 4, Cambridge University Press, 2006. \href{https://doi.org/10.1017/CBO9780511542879}{doi:10.1017/CBO9780511542879}
\bibitem[BK10]{BannaiKobayashi2010} \textsc{Kenichi Bannai and Shinichi Kobayashi}, Algebraic theta functions and the $p$-adic interpolation of Eisenstein--Kronecker numbers, \emph{Duke Math. J.} \textbf{153} (2010), 229--295. \href{https://doi.org/10.1215/00127094-2010-024}{doi:10.1215/00127094-2010-024}
\bibitem[BK90]{BlochKato1990} \textsc{Spencer Bloch and Kazuya Kato}, $L$-functions and Tamagawa numbers of motives, in \emph{The Grothendieck Festschrift, Vol. I}, Progr. Math. 86, Birkh\"{a}user, 1990, 333--400.
\bibitem[BLGGT14]{BLGGT2014} \textsc{Thomas Barnet-Lamb, Toby Gee, David Geraghty, and Richard Taylor}, Local-global compatibility for $l=p$, II, \emph{Ann. Sci. \'{E}c. Norm. Sup\'{e}r. (4)} \textbf{47} (2014), 165--179. \href{https://doi.org/10.24033/asens.2212}{doi:10.24033/asens.2212}
\bibitem[BMS18]{BhattMorrowScholze2018} \textsc{Bhargav Bhatt, Matthew Morrow, and Peter Scholze}, Integral $p$-adic Hodge theory, \emph{Publ. Math. Inst. Hautes \'{E}tudes Sci.} \textbf{128} (2018), 219--397. \href{https://doi.org/10.1007/s10240-019-00102-z}{doi:10.1007/s10240-019-00102-z}
\bibitem[Bom00]{Bombieri2000} \textsc{Enrico Bombieri}, Problems of the Millennium: the Riemann Hypothesis, Clay Mathematics Institute, 2000. \url{https://www.claymath.org/wp-content/uploads/2022/05/riemann.pdf}
\bibitem[Bom74]{Bombieri1974} \textsc{Enrico Bombieri}, Counting points on curves over finite fields (d'apr\`{e}s S. A. Stepanov), in \emph{S\'{e}minaire Bourbaki, vol. 1972/73, expos\'{e} 430}, Lecture Notes in Math. 383, Springer, 1974, 234--241. \href{https://doi.org/10.1007/BFb0057311}{doi:10.1007/BFb0057311}
\bibitem[Bor07]{Borel2007} \textsc{Armand Borel}, Automorphic Forms on Reductive Groups, in \emph{Automorphic Forms and Applications}, IAS/Park City Math. Ser. 12, American Mathematical Society, 2007, 5--39. \href{https://doi.org/10.1090/pcms/012/02}{doi:10.1090/pcms/012/02}
\bibitem[Bos01]{Bost2001} \textsc{Jean-Beno\^{\i}t Bost}, Algebraic leaves of algebraic foliations over number fields, \emph{Publ. Math. Inst. Hautes \'{E}tudes Sci.} \textbf{93} (2001), 161--221. \href{https://doi.org/10.1007/s10240-001-8191-3}{doi:10.1007/s10240-001-8191-3}
\bibitem[Bra41]{Braun1941} \textsc{Hel Braun}, Zur Theorie der hermitischen Formen, \emph{Abh. Math. Sem. Univ. Hamburg} \textbf{14} (1941), 61--150. \href{https://doi.org/10.1007/BF02940742}{doi:10.1007/BF02940742}
\bibitem[BV83]{BombieriVaaler1983} \textsc{Enrico Bombieri and Jeffrey D. Vaaler}, On Siegel's lemma, \emph{Invent. Math.} \textbf{73} (1983), 11--32. \href{https://doi.org/10.1007/BF01393823}{doi:10.1007/BF01393823}
\bibitem[CH13]{ChenevierHarris2013} \textsc{Ga\"{e}tan Chenevier and Michael Harris}, Construction of automorphic Galois representations, II, \emph{Camb. J. Math.} \textbf{1} (2013), 53--73. \href{https://doi.org/10.4310/CJM.2013.v1.n1.a2}{doi:10.4310/CJM.2013.v1.n1.a2}
\bibitem[Col98]{Colmez1998} \textsc{Pierre Colmez}, Th\'{e}orie d'Iwasawa des repr\'{e}sentations de de Rham d'un corps local, \emph{Ann. of Math. (2)} \textbf{148} (1998), 485--571. \href{https://doi.org/10.2307/121003}{doi:10.2307/121003}
\bibitem[Con03]{Conrey2003} \textsc{J. Brian Conrey}, The Riemann hypothesis, \emph{Notices Amer. Math. Soc.} \textbf{50} (2003), 341--353. \url{https://www.ams.org/notices/200303/fea-conrey-web.pdf}
\bibitem[Con89]{Conrey1989} \textsc{J. Brian Conrey}, More than two fifths of the zeros of the Riemann zeta function are on the critical line, \emph{J. Reine Angew. Math.} \textbf{399} (1989), 1--26. \href{https://doi.org/10.1515/crll.1989.399.1}{doi:10.1515/crll.1989.399.1}
\bibitem[Dav80]{Davenport1980} \textsc{Harold Davenport}, \emph{Multiplicative Number Theory}, 2nd ed., Graduate Texts in Math. 74, Springer, 1980, revised by H. L. Montgomery. \href{https://doi.org/10.1007/978-1-4757-5927-3}{doi:10.1007/978-1-4757-5927-3}
\bibitem[Del71]{Deligne1971} \textsc{Pierre Deligne}, Formes modulaires et repr\'{e}sentations $\ell$-adiques, in \emph{S\'{e}minaire Bourbaki 1968/69, Exp. 355}, Lecture Notes in Math. 179, Springer, 1971, 139--172.
\bibitem[Del74]{Deligne1974} \textsc{Pierre Deligne}, La conjecture de Weil. I, \emph{Publ. Math. Inst. Hautes \'{E}tudes Sci.} \textbf{43} (1974), 273--307. \href{https://doi.org/10.1007/BF02684373}{doi:10.1007/BF02684373}
\bibitem[Del80]{Deligne1980} \textsc{Pierre Deligne}, La conjecture de Weil. II, \emph{Publ. Math. Inst. Hautes \'{E}tudes Sci.} \textbf{52} (1980), 137--252. \href{https://doi.org/10.1007/BF02684780}{doi:10.1007/BF02684780}
\bibitem[dlVP96]{ValleePoussin1896} \textsc{Charles-Jean de la Vall\'{e}e Poussin}, Recherches analytiques sur la th\'{e}orie des nombres premiers, \emph{Ann. Soc. Sci. Bruxelles} \textbf{20} (1896), 183--256, 281--397.
\bibitem[dMU21]{deMouraUllrich2021} \textsc{Leonardo de Moura and Sebastian Ullrich}, The Lean 4 theorem prover and programming language, in \emph{Automated Deduction -- CADE 28}, Lecture Notes in Comput. Sci. 12699, Springer, 2021, 625--635. \href{https://doi.org/10.1007/978-3-030-79876-5_37}{doi:10.1007/978-3-030-79876-5\_37}
\bibitem[Edw74]{Edwards1974} \textsc{Harold M. Edwards}, \emph{Riemann's Zeta Function}, Pure and Applied Mathematics 58, Academic Press, 1974.
\bibitem[FJ80]{FlenstedJensen1980} \textsc{Mogens Flensted-Jensen}, Discrete series for semisimple symmetric spaces, \emph{Ann. of Math. (2)} \textbf{111} (1980), 253--311. \href{https://doi.org/10.2307/1971201}{doi:10.2307/1971201}
\bibitem[Ful97]{Fulton1997} \textsc{William Fulton}, \emph{Young Tableaux}, London Math. Soc. Stud. Texts 35, Cambridge University Press, 1997. \href{https://doi.org/10.1017/CBO9780511626241}{doi:10.1017/CBO9780511626241}
\bibitem[FZ21]{Forrester2021} \textsc{Peter J. Forrester and Jiyuan Zhang}, Corank-1 projections and the randomised Horn problem, \emph{Tunis. J. Math.} \textbf{3} (2021), 55--73. \href{https://doi.org/10.2140/tunis.2021.3.55}{doi:10.2140/tunis.2021.3.55}
\bibitem[Gau08]{Gaudron2008} \textsc{\'{E}ric Gaudron}, Pentes des fibr\'{e}s vectoriels ad\'{e}liques sur un corps global, \emph{Rend. Sem. Mat. Univ. Padova} \textbf{119} (2008), 21--95. \href{https://doi.org/10.4171/RSMUP/119-2}{doi:10.4171/RSMUP/119-2}
\bibitem[GGP12]{GGP2012} \textsc{Wee Teck Gan, Benedict H. Gross, and Dipendra Prasad}, Symplectic local root numbers, central critical $L$-values, and restriction problems in the representation theory of classical groups, in \emph{Sur les conjectures de Gross et Prasad. I}, Ast\'{e}risque 346, Soc. Math. France, 2012, 1--109.
\bibitem[GHY01]{GanHankeYu2001} \textsc{Wee Teck Gan, Jonathan P. Hanke, and Jiu-Kang Yu}, On an exact mass formula of Shimura, \emph{Duke Math. J.} \textbf{107} (2001), 103--133. \href{https://doi.org/10.1215/S0012-7094-01-10716-3}{doi:10.1215/S0012-7094-01-10716-3}
\bibitem[GI16]{GanIchino2016} \textsc{Wee Teck Gan and Atsushi Ichino}, The Gross--Prasad conjecture and local theta correspondence, \emph{Invent. Math.} \textbf{206} (2016), 705--799. \href{https://doi.org/10.1007/s00222-016-0662-8}{doi:10.1007/s00222-016-0662-8}
\bibitem[Gre55]{Green1955} \textsc{James A. Green}, The characters of the finite general linear groups, \emph{Trans. Amer. Math. Soc.} \textbf{80} (1955), 402--447. \href{https://doi.org/10.1090/S0002-9947-1955-0072878-2}{doi:10.1090/S0002-9947-1955-0072878-2}
\bibitem[Gre80]{Green1980} \textsc{James A. Green}, \emph{Polynomial Representations of $\mathrm{GL}_n$}, Lecture Notes in Math. 830, Springer, 1980. \href{https://doi.org/10.1007/BFb0092296}{doi:10.1007/BFb0092296}
\bibitem[Gre89]{Greenberg1989} \textsc{Ralph Greenberg}, Iwasawa theory for $p$-adic representations, in \emph{Algebraic Number Theory}, Adv. Stud. Pure Math. 17, Academic Press, 1989, 97--137.
\bibitem[Gro99]{Gross1999} \textsc{Benedict H. Gross}, Algebraic modular forms, \emph{Israel J. Math.} \textbf{113} (1999), 61--93. \href{https://doi.org/10.1007/BF02780173}{doi:10.1007/BF02780173}
\bibitem[GW09]{GoodmanWallach2009} \textsc{Roe Goodman and Nolan R. Wallach}, \emph{Symmetry, Representations, and Invariants}, Grad. Texts in Math. 255, Springer, 2009. \href{https://doi.org/10.1007/978-0-387-79852-3}{doi:10.1007/978-0-387-79852-3}
\bibitem[Had96]{Hadamard1896} \textsc{Jacques Hadamard}, Sur la distribution des z\'{e}ros de la fonction $\zeta(s)$ et ses cons\'{e}quences arithm\'{e}tiques, \emph{Bull. Soc. Math. France} \textbf{24} (1896), 199--220. \href{https://doi.org/10.24033/bsmf.545}{doi:10.24033/bsmf.545}
\bibitem[Har14]{Hardy1914} \textsc{Godfrey Harold Hardy}, Sur les z\'{e}ros de la fonction $\zeta(s)$ de Riemann, \emph{C. R. Acad. Sci. Paris} \textbf{158} (1914), 1012--1014.
\bibitem[HB95]{HeathBrown1995} \textsc{D. Roger Heath-Brown}, A mean value estimate for real character sums, \emph{Acta Arith.} \textbf{72} (1995), 235--275. \href{https://doi.org/10.4064/aa-72-3-235-275}{doi:10.4064/aa-72-3-235-275}
\bibitem[He17]{He2017} \textsc{Hongyu He}, On the Gan--Gross--Prasad conjecture for $U(p,q)$, \emph{Invent. Math.} \textbf{209} (2017), 837--884. \href{https://doi.org/10.1007/s00222-017-0720-x}{doi:10.1007/s00222-017-0720-x}
\bibitem[HL04]{HarrisLabesse2004} \textsc{Michael Harris and Jean-Pierre Labesse}, Conditional base change for unitary groups, \emph{Asian J. Math.} \textbf{8} (2004), 653--684. \href{https://doi.org/10.4310/AJM.2004.v8.n4.a19}{doi:10.4310/AJM.2004.v8.n4.a19}
\bibitem[How89]{Howe1989} \textsc{Roger Howe}, Transcending classical invariant theory, \emph{J. Amer. Math. Soc.} \textbf{2} (1989), 535--552. \href{https://doi.org/10.1090/S0894-0347-1989-0985172-6}{doi:10.1090/S0894-0347-1989-0985172-6}
\bibitem[HT01]{HarrisTaylor2001} \textsc{Michael Harris and Richard Taylor}, \emph{The Geometry and Cohomology of Some Simple Shimura Varieties}, Ann. of Math. Stud. 151, Princeton University Press, 2001. \href{https://doi.org/10.1515/9781400837205}{doi:10.1515/9781400837205}
\bibitem[Hua63]{Hua1963} \textsc{Loo-Keng Hua}, \emph{Harmonic Analysis of Functions of Several Complex Variables in the Classical Domains}, Transl. Math. Monogr. 6, American Mathematical Society, 1963. \href{https://doi.org/10.1090/mmono/006}{doi:10.1090/mmono/006}
\bibitem[IK04]{IwaniecKowalski2004} \textsc{Henryk Iwaniec and Emmanuel Kowalski}, \emph{Analytic Number Theory}, Amer. Math. Soc. Colloq. Publ. 53, American Mathematical Society, 2004. \href{https://doi.org/10.1090/coll/053}{doi:10.1090/coll/053}
\bibitem[Jut75]{Jutila1975} \textsc{Matti Jutila}, On mean values of Dirichlet polynomials with real characters, \emph{Acta Arith.} \textbf{27} (1975), 191--198. \href{https://doi.org/10.4064/aa-27-1-191-198}{doi:10.4064/aa-27-1-191-198}
\bibitem[Kis06]{Kisin2006} \textsc{Mark Kisin}, Crystalline representations and $F$-crystals, in \emph{Algebraic Geometry and Number Theory}, Progr. Math. 253, Birkh\"{a}user, 2006, 459--496.
\bibitem[KL81]{KubertLang1981} \textsc{Daniel S. Kubert and Serge Lang}, \emph{Modular Units}, Grundlehren Math. Wiss. 244, Springer, 1981. \href{https://doi.org/10.1007/978-1-4757-1741-9}{doi:10.1007/978-1-4757-1741-9}
\bibitem[Klo09]{Klosin2009} \textsc{Krzysztof Klosin}, Congruences among modular forms on $U(2,2)$ and the Bloch--Kato conjecture, \emph{Ann. Inst. Fourier (Grenoble)} \textbf{59} (2009), 81--166. \href{https://doi.org/10.5802/aif.2427}{doi:10.5802/aif.2427}
\bibitem[Kne65]{Kneser1965} \textsc{Martin Kneser}, Starke Approximation in algebraischen Gruppen. I, \emph{J. Reine Angew. Math.} \textbf{218} (1965), 190--203. \href{https://doi.org/10.1515/crll.1965.218.190}{doi:10.1515/crll.1965.218.190}
\bibitem[Kot92]{Kottwitz1992} \textsc{Robert E. Kottwitz}, Points on some Shimura varieties over finite fields, \emph{J. Amer. Math. Soc.} \textbf{5} (1992), 373--444. \href{https://doi.org/10.1090/S0894-0347-1992-1124982-1}{doi:10.1090/S0894-0347-1992-1124982-1}
\bibitem[Kud84]{Kudla1984} \textsc{Stephen S. Kudla}, Seesaw dual reductive pairs, in \emph{Automorphic Forms of Several Variables (Katata, 1983)}, Progr. Math. 46, Birkh\"{a}user, 1984, 244--268.
\bibitem[Kud86]{Kudla1986} \textsc{Stephen S. Kudla}, On the local theta-correspondence, \emph{Invent. Math.} \textbf{83} (1986), 229--255. \href{https://doi.org/10.1007/BF01388961}{doi:10.1007/BF01388961}
\bibitem[KV78]{KashiwaraVergne1978} \textsc{Masaki Kashiwara and Mich\`{e}le Vergne}, On the Segal--Shale--Weil representations and harmonic polynomials, \emph{Invent. Math.} \textbf{44} (1978), 1--47. \href{https://doi.org/10.1007/BF01389900}{doi:10.1007/BF01389900}
\bibitem[Lan05]{Landau1905} \textsc{Edmund Landau}, \"{U}ber einen Satz von Tschebyschef, \emph{Math. Ann.} \textbf{61} (1905), 527--550. \href{https://doi.org/10.1007/BF01449495}{doi:10.1007/BF01449495}
\bibitem[Lau94]{Laurent1994} \textsc{Michel Laurent}, Linear forms in two logarithms and interpolation determinants, \emph{Acta Arith.} \textbf{66} (1994), 181--199. \href{https://doi.org/10.4064/aa-66-2-181-199}{doi:10.4064/aa-66-2-181-199}
\bibitem[Lev74]{Levinson1974} \textsc{Norman Levinson}, More than one third of zeros of Riemann's zeta-function are on $\sigma=1/2$, \emph{Adv. Math.} \textbf{13} (1974), 383--436. \href{https://doi.org/10.1016/0001-8708(74)90074-7}{doi:10.1016/0001-8708(74)90074-7}
\bibitem[Li90]{Li1990} \textsc{Jian-Shu Li}, Theta lifting for unitary representations with nonzero cohomology, \emph{Duke Math. J.} \textbf{61} (1990), 913--937. \href{https://doi.org/10.1215/S0012-7094-90-06135-6}{doi:10.1215/S0012-7094-90-06135-6}
\bibitem[Li92]{Li1992} \textsc{Jian-Shu Li}, Non-vanishing theorems for the cohomology of certain arithmetic quotients, \emph{J. Reine Angew. Math.} \textbf{428} (1992), 177--217. \href{https://doi.org/10.1515/crll.1992.428.177}{doi:10.1515/crll.1992.428.177}
\bibitem[Liu26]{S03} \textsc{Vincent Liu}, \emph{Arithmetic Spectral Descent: Obstruction Theory, Hodge--Krein Realization, and a Falsifiable Architecture for the Riemann Hypothesis}, Zenodo, 2026, WiseMotif Monograph of Mathematical Intelligence, Series 03; 4 volumes. \href{https://doi.org/10.5281/zenodo.23564898}{doi:10.5281/zenodo.23564898}
\bibitem[LT65]{LubinTate1965} \textsc{Jonathan Lubin and John Tate}, Formal complex multiplication in local fields, \emph{Ann. of Math. (2)} \textbf{81} (1965), 380--387. \href{https://doi.org/10.2307/1970622}{doi:10.2307/1970622}
\bibitem[Man95]{vonMangoldt1895} \textsc{Hans von Mangoldt}, Zu Riemanns Abhandlung \quotedblbase{}Ueber die Anzahl der Primzahlen unter einer gegebenen Gr\"{o}sse``, \emph{J. Reine Angew. Math.} \textbf{114} (1895), 255--305. \href{https://doi.org/10.1515/crll.1895.114.255}{doi:10.1515/crll.1895.114.255}
\bibitem[mC20]{mathlib2020} \textsc{The mathlib Community}, The Lean mathematical library, in \emph{Proceedings of the 9th ACM SIGPLAN International Conference on Certified Programs and Proofs (CPP 2020)}, ACM, 2020, 367--381. \href{https://doi.org/10.1145/3372885.3373824}{doi:10.1145/3372885.3373824}
\bibitem[Mil06]{Milne2006} \textsc{James S. Milne}, \emph{Arithmetic Duality Theorems}, 2nd ed., BookSurge, 2006.
\bibitem[MR10]{MazurRubin2010} \textsc{Barry Mazur and Karl Rubin}, Ranks of twists of elliptic curves and Hilbert's tenth problem, \emph{Invent. Math.} \textbf{181} (2010), 541--575. \href{https://doi.org/10.1007/s00222-010-0252-0}{doi:10.1007/s00222-010-0252-0}
\bibitem[MS12]{MillerSchmid2012} \textsc{Stephen D. Miller and Wilfried Schmid}, On the rapid decay of cuspidal automorphic forms, \emph{Adv. Math.} \textbf{231} (2012), 940--964. \href{https://doi.org/10.1016/j.aim.2012.06.001}{doi:10.1016/j.aim.2012.06.001}
\bibitem[Mum83]{Mumford1983} \textsc{David Mumford}, \emph{Tata Lectures on Theta I}, Progr. Math. 28, Birkh\"{a}user, 1983. \href{https://doi.org/10.1007/978-1-4899-2843-6}{doi:10.1007/978-1-4899-2843-6}
\bibitem[MV07]{MontgomeryVaughan2007} \textsc{Hugh L. Montgomery and Robert C. Vaughan}, \emph{Multiplicative Number Theory I. Classical Theory}, Cambridge Stud. Adv. Math. 97, Cambridge University Press, 2007. \href{https://doi.org/10.1017/CBO9780511618314}{doi:10.1017/CBO9780511618314}
\bibitem[NV21]{NelsonVenkatesh2021} \textsc{Paul D. Nelson and Akshay Venkatesh}, The orbit method and analysis of automorphic forms, \emph{Acta Math.} \textbf{226} (2021), 1--209. \href{https://doi.org/10.4310/ACTA.2021.v226.n1.a1}{doi:10.4310/ACTA.2021.v226.n1.a1}
\bibitem[O'M73]{OMeara1973} \textsc{O. Timothy O'Meara}, \emph{Introduction to Quadratic Forms}, 3rd corrected printing ed., Grundlehren Math. Wiss. 117, Springer, 1973. \href{https://doi.org/10.1007/978-3-662-41922-9}{doi:10.1007/978-3-662-41922-9}
\bibitem[OT51]{OstrowskiTaussky1951} \textsc{Alexander M. Ostrowski and Olga Taussky}, On the variation of the determinant of a positive definite matrix, \emph{Nederl. Akad. Wetensch. Proc. Ser. A} \textbf{54} (1951), 383--385.
\bibitem[Pau99]{Paul1999} \textsc{Annegret Paul}, First occurrence for the dual pairs $(U(p,q),U(r,s))$, \emph{Canad. J. Math.} \textbf{51} (1999), 636--657. \href{https://doi.org/10.4153/CJM-1999-029-6}{doi:10.4153/CJM-1999-029-6}
\bibitem[Pra06]{Prasanna2006} \textsc{Kartik Prasanna}, Integrality of a ratio of Petersson norms and level-lowering congruences, \emph{Ann. of Math. (2)} \textbf{163} (2006), 901--967. \href{https://doi.org/10.4007/annals.2006.163.901}{doi:10.4007/annals.2006.163.901}
\bibitem[PT21]{PlattTrudgian2021} \textsc{David J. Platt and Timothy S. Trudgian}, The Riemann hypothesis is true up to $3\cdot10^{12}$, \emph{Bull. Lond. Math. Soc.} \textbf{53} (2021), 792--797. \href{https://doi.org/10.1112/blms.12460}{doi:10.1112/blms.12460}
\bibitem[Ral82]{Rallis1982} \textsc{Stephen Rallis}, Langlands' functoriality and the Weil representation, \emph{Amer. J. Math.} \textbf{104} (1982), 469. \href{https://doi.org/10.2307/2374151}{doi:10.2307/2374151}
\bibitem[Ral84]{Rallis1984} \textsc{Stephen Rallis}, Injectivity properties of liftings associated to Weil representations, \emph{Compositio Math.} \textbf{52} (1984), 139--169.
\bibitem[Ral84b]{RallisHowe1984} \textsc{Stephen Rallis}, On the Howe duality conjecture, \emph{Compositio Math.} \textbf{51} (1984), 333--399.
\bibitem[Ral87]{Rallis1987} \textsc{Stephen Rallis}, \emph{L-Functions and the Oscillator Representation}, Lecture Notes in Math. 1245, Springer, 1987. \href{https://doi.org/10.1007/BFb0077894}{doi:10.1007/BFb0077894}
\bibitem[Rib76]{Ribet1976} \textsc{Kenneth A. Ribet}, A modular construction of unramified $p$-extensions of $\mathbf Q(\mu_p)$, \emph{Invent. Math.} \textbf{34} (1976), 151--162. \href{https://doi.org/10.1007/BF01403065}{doi:10.1007/BF01403065}
\bibitem[Rie59]{Riemann1859} \textsc{Bernhard Riemann}, Ueber die Anzahl der Primzahlen unter einer gegebenen Gr\"{o}sse, \emph{Monatsber. K\"{o}nigl. Preuss. Akad. Wiss. Berlin} (1859), 671--680. \url{https://www.claymath.org/collections/riemanns-1859-manuscript/}
\bibitem[Rou96]{Rouquier1996} \textsc{Rapha\"{e}l Rouquier}, Caract\'{e}risation des caract\`{e}res et pseudo-caract\`{e}res, \emph{J. Algebra} \textbf{180} (1996), 571--586. \href{https://doi.org/10.1006/jabr.1996.0083}{doi:10.1006/jabr.1996.0083}
\bibitem[SC67]{ChowlaSelberg1967} \textsc{Atle Selberg and Sarvadaman Chowla}, On Epstein's zeta-function, \emph{J. Reine Angew. Math.} \textbf{227} (1967), 86--110. \href{https://doi.org/10.1515/crll.1967.227.86}{doi:10.1515/crll.1967.227.86}
\bibitem[Sch13]{Scholze2013} \textsc{Peter Scholze}, $p$-adic Hodge theory for rigid-analytic varieties, \emph{Forum Math. Pi} \textbf{1} (2013), e1. \href{https://doi.org/10.1017/fmp.2013.1}{doi:10.1017/fmp.2013.1}
\bibitem[Sel42]{Selberg1942} \textsc{Atle Selberg}, On the zeros of Riemann's zeta-function, \emph{Skr. Norske Vid.-Akad. Oslo I} \textbf{1942} (1942), no. 10, 1--59.
\bibitem[Ser97]{Serre1997} \textsc{Jean-Pierre Serre}, \emph{Galois Cohomology}, Springer Monogr. Math., Springer, 1997. \href{https://doi.org/10.1007/978-3-642-59141-9}{doi:10.1007/978-3-642-59141-9}
\bibitem[Shi11]{Shin2011} \textsc{Sug Woo Shin}, Galois representations arising from some compact Shimura varieties, \emph{Ann. of Math. (2)} \textbf{173} (2011), 1645--1741. \href{https://doi.org/10.4007/annals.2011.173.3.9}{doi:10.4007/annals.2011.173.3.9}
\bibitem[Shi71]{Shimura1971} \textsc{Goro Shimura}, \emph{Introduction to the Arithmetic Theory of Automorphic Functions}, Publ. Math. Soc. Japan 11, Iwanami Shoten and Princeton University Press, 1971.
\bibitem[Shi76]{Shimura1976} \textsc{Goro Shimura}, Theta functions with complex multiplication, \emph{Duke Math. J.} \textbf{43} (1976), 673--696. \href{https://doi.org/10.1215/S0012-7094-76-04353-2}{doi:10.1215/S0012-7094-76-04353-2}
\bibitem[Sie35]{Siegel1935} \textsc{Carl Ludwig Siegel}, \"{U}ber die analytische Theorie der quadratischen Formen, \emph{Ann. of Math. (2)} \textbf{36} (1935), 527--606. \href{https://doi.org/10.2307/1968644}{doi:10.2307/1968644}
\bibitem[Sil09]{Silverman2009} \textsc{Joseph H. Silverman}, \emph{The Arithmetic of Elliptic Curves}, 2nd ed., Graduate Texts in Math. 106, Springer, 2009. \href{https://doi.org/10.1007/978-0-387-09494-6}{doi:10.1007/978-0-387-09494-6}
\bibitem[Sil94]{Silverman1994} \textsc{Joseph H. Silverman}, \emph{Advanced Topics in the Arithmetic of Elliptic Curves}, Graduate Texts in Math. 151, Springer, 1994. \href{https://doi.org/10.1007/978-1-4612-0851-8}{doi:10.1007/978-1-4612-0851-8}
\bibitem[Ste69]{Stepanov1969} \textsc{Sergei A. Stepanov}, On the number of points of a hyperelliptic curve over a finite prime field, \emph{Math. USSR-Izv.} \textbf{3} (1969), 1103--1114, translated from Izv. Akad. Nauk SSSR Ser. Mat. 33 (1969), 1171--1181. \href{https://doi.org/10.1070/IM1969v003n05ABEH000834}{doi:10.1070/IM1969v003n05ABEH000834}
\bibitem[Sun07]{Sun2007} \textsc{Binyong Sun}, Matrix coefficients of cohomologically induced representations, \emph{Compos. Math.} \textbf{143} (2007), 201--221. \href{https://doi.org/10.1112/S0010437X06002508}{doi:10.1112/S0010437X06002508}
\bibitem[Suz24]{Suzuki2024} \textsc{Masatoshi Suzuki}, On variants of Chebyshev's conjecture, preprint, 2024. \href{https://arxiv.org/abs/2411.07436}{arXiv:2411.07436}
\bibitem[Tad86]{Tadic1986} \textsc{Marko Tadi\'{c}}, Classification of unitary representations in irreducible representations of general linear group (non-Archimedean case), \emph{Ann. Sci. \'{E}cole Norm. Sup. (4)} \textbf{19} (1986), 335--382. \href{https://doi.org/10.24033/asens.1510}{doi:10.24033/asens.1510}
\bibitem[Tat67]{Tate1967} \textsc{John T. Tate}, Fourier analysis in number fields, and Hecke's zeta-functions, in \emph{Algebraic Number Theory (J. W. S. Cassels and A. Fr\"{o}hlich, eds.)}, Academic Press, 1967, 305--347, thesis, Princeton, 1950.
\bibitem[Tat67b]{TatePDiv1967} \textsc{John T. Tate}, $p$-divisible groups, in \emph{Proceedings of a Conference on Local Fields (Driebergen, 1966)}, Springer, 1967, 158--183. \href{https://doi.org/10.1007/978-3-642-87942-5_12}{doi:10.1007/978-3-642-87942-5\_12}
\bibitem[Tay91]{Taylor1991} \textsc{Richard Taylor}, Galois representations associated to Siegel modular forms of low weight, \emph{Duke Math. J.} \textbf{63} (1991), 281--332. \href{https://doi.org/10.1215/S0012-7094-91-06312-X}{doi:10.1215/S0012-7094-91-06312-X}
\bibitem[Tit86]{Titchmarsh1986} \textsc{Edward Charles Titchmarsh}, \emph{The Theory of the Riemann Zeta-Function}, 2nd ed., Oxford University Press, 1986, revised by D. R. Heath-Brown.
\bibitem[Wac96]{Wach1996} \textsc{Nathalie Wach}, Repr\'{e}sentations $p$-adiques potentiellement cristallines, \emph{Bull. Soc. Math. France} \textbf{124} (1996), 375--400. \href{https://doi.org/10.24033/bsmf.2285}{doi:10.24033/bsmf.2285}
\bibitem[Wal15]{Walls2015} \textsc{Patrick Walls}, The theta correspondence, periods of automorphic forms and special values of standard automorphic $L$-functions, preprint, 2015. \href{https://arxiv.org/abs/1501.01565}{arXiv:1501.01565}
\bibitem[Wal88]{Wallach1988} \textsc{Nolan R. Wallach}, \emph{Real Reductive Groups I}, Pure and Applied Mathematics 132, Academic Press, 1988.
\bibitem[Wei48]{Weil1948} \textsc{Andr\'{e} Weil}, \emph{Sur les courbes alg\'{e}briques et les vari\'{e}t\'{e}s qui s'en d\'{e}duisent}, Actualit\'{e}s Sci. Ind. 1041, Hermann, 1948.
\bibitem[Wei49]{Weil1949} \textsc{Andr\'{e} Weil}, Numbers of solutions of equations in finite fields, \emph{Bull. Amer. Math. Soc.} \textbf{55} (1949), 497--508. \href{https://doi.org/10.1090/S0002-9904-1949-09219-4}{doi:10.1090/S0002-9904-1949-09219-4}
\bibitem[Wei52]{Weil1952} \textsc{Andr\'{e} Weil}, Sur les \guillemotleft{} formules explicites \guillemotright{} de la th\'{e}orie des nombres premiers, \emph{Comm. S\'{e}m. Math. Univ. Lund, tome suppl\'{e}mentaire} (1952), 252--265.
\bibitem[Wei56]{Weil1956} \textsc{Andr\'{e} Weil}, On a certain type of characters of the id\`{e}le-class group of an algebraic number-field, in \emph{Proceedings of the International Symposium on Algebraic Number Theory, Tokyo--Nikko 1955}, Science Council of Japan, 1956, 1--7.
\bibitem[Wei64]{Weil1964} \textsc{Andr\'{e} Weil}, Sur certains groupes d'op\'{e}rateurs unitaires, \emph{Acta Math.} \textbf{111} (1964), 143--211. \href{https://doi.org/10.1007/BF02391012}{doi:10.1007/BF02391012}
\bibitem[Wei65]{Weil1965} \textsc{Andr\'{e} Weil}, Sur la formule de Siegel dans la th\'{e}orie des groupes classiques, \emph{Acta Math.} \textbf{113} (1965), 1--87. \href{https://doi.org/10.1007/BF02391774}{doi:10.1007/BF02391774}
\bibitem[Wid41]{Widder1941} \textsc{David Vernon Widder}, \emph{The Laplace Transform}, Princeton Math. Ser. 6, Princeton University Press, 1941.
\bibitem[Wil90]{Wiles1990} \textsc{Andrew Wiles}, The Iwasawa conjecture for totally real fields, \emph{Ann. of Math. (2)} \textbf{131} (1990), 493--540. \href{https://doi.org/10.2307/1971468}{doi:10.2307/1971468}
\bibitem[Wol82]{Wolsey1982} \textsc{Laurence A. Wolsey}, An analysis of the greedy algorithm for the submodular set covering problem, \emph{Combinatorica} \textbf{2} (1982), 385--393. \href{https://doi.org/10.1007/BF02579435}{doi:10.1007/BF02579435}
\bibitem[Wu13]{Wu2013} \textsc{Chenyan Wu}, Irreducibility of theta lifting for unitary groups, \emph{J. Number Theory} \textbf{133} (2013), 3296--3318. \href{https://doi.org/10.1016/j.jnt.2013.04.014}{doi:10.1016/j.jnt.2013.04.014}
\end{thebibliography}
